\begin{table}[t]
\centering
\small
\begin{tabular}{lcccccc}
\toprule
Method & Exact & Contam. & Trunc. & Over-exp. & Suffix/Num. & Drift \\
\midrule
\textit{Frontier DeepSeek} & 12 & 0 & 1 & 1 & 9 & 27 \\
MT-SFT & 5 & 1 & 0 & 5 & 12 & 27 \\
\textit{MT-DPO g0.4 sig} & 5 & 2 & 0 & 2 & 7 & 34 \\
\textit{Open selector (3-way, open)} & 12 & 0 & 0 & 1 & 8 & 29 \\
Catalog adjudicator (3-way, closed) & 14 & 0 & 0 & 0 & 7 & 29 \\
Legacy DPO & 2 & 9 & 0 & 2 & 8 & 29 \\
\bottomrule
\end{tabular}
\caption{Deterministic failure taxonomy on the 50-item Tangut test set. ``Contam.'' is a binary artifact flag with no tuned threshold: any Latin alphabetic character, angle-bracket markup, explicit \texttt{[UNK]}, or one of the small prompt-artifact lexemes observed in legacy runs marks the output as contaminated. ``Suffix/Num.'' groups still-title-like cases that either miss at least one reference-side title suffix or mismatch the reference numeral string. ``Drift'' denotes the remaining non-exact cases after these categories are removed. The table makes the failure modes more specific than the scalar contamination rate alone: frontier prompting is already clean, legacy DPO is uniquely artifact-prone, and the open selector keeps the clean frontier regime while removing the frontier truncation case.}
\label{tab:failure-taxonomy}
\end{table}
